\section{Ateji　当て字}

What is the Ateji? It is a pan-sinosphere way of using Chinese characters to represent sounds of a word, where the Chinese characters chosen to represent the sound are themselves semantically related to the meaning of the word itself. It is a usage that is seen in every sinoglyph-using writing system, and is intimately related to the evolution of the phonosemantic character, and the Kana syllabary. The Japanese term of the phenonomenon is chosen to ground our discussion because the Japanese have the most evolved, sophisticated, and widespread usage of ateji in the Sinosphere.  

In typical usage, however, the term "ateji" refers to the Japanese practice of using Chinese characters for their sound value, regardless of whether the characters chosen are semantically related to the meaning of the word itself. In some cases, the sinoglyphs chosen do indeed semantically hint or are related to the meaning of the word itself. In our discussion, we will only consider the more restricted case - where the sinoglyphs chosen are indeed semantically related to the meaning of the word itself - and we will treat "ateji" to mean that.

Here are some examples of ateji in Japanese 
護美　ゴミ rubbish
倶楽部　クラブ club

In the manyoshu 
烏梅能波奈

0816　【承前，卅二之二。】
　烏梅能波奈　伊麻佐家留期等　知利須義受　和我霸能曾能爾　阿利己世奴加毛　【少貳小野大夫。】
　梅花うめのはな　今咲いまさける如ごと　散過ちりすぎず　我わが家園へのそのに　在ありこせぬかも　【少貳小野大夫せうにをのだいぶ。】
　今見此梅花　盛咲之後未散盡　風情洽宜矣　還願吾家庭園間　殘梅如是莫盡落

珍紛漢紛ちんぷんかんぷん・珍糞漢糞ちんぷんかんぷん・陳奮翰奮（ちんぷんかんぷん）[6]
出鱈目（でたらめ）[6]
時計（とけい）[6]
馬鹿（ばか）[6]
滅茶苦茶（めちゃくちゃ）[6]
滅多（めった）[6]

また、夏目漱石が、当て字をよく用いていた作家として有名である。

瓦楽多（がらくた）[9]
伽藍堂（がらんどう）[9]
愚図愚図（ぐずぐず）[9]
頓痴気（とんちき）[9]
浪漫（ロマン）[10]


大丈夫＝だいじょぶ

Mandarin examples:
威而剛 = Viagra
万维网 = World Wide Web 
黑客 = Hacker

雷射 = Laser
可口可樂

彭定康 = Christopher Patten

世界語 = 愛斯不難讀 = Esperanto

英吉利 德意志 美利堅合眾國
Shanghainese
江摆渡（kaonpadu / comprador）-洋行買辦江摆渡（kaonpadu / comprador），小輪叫司汀巴（sythinpo / steamer）。

阿斯匹靈（aspirin）

德律風（telephone）

馬達（motor）

幽默（humour）
\begin{definition}[Ateji]
    An object $s \in S$ is an ateji if $s$ is a emergent object $s= a_{x}$, i.e. $s\in\langle a \rangle$, where $x = x_{1} * \cdots * x_{n}$ (i.e. $x$ is a composite object)and $s= a_{x} \sim x_{1} * \cdots * x_{n}$.

    In other words, an object $s$ is an ateji if and only if it is an emergent object, and is synonymous with a composite object.

    An symbol $s \in S$ is an ateji if $s$ is a phonosemantic character $s= a_{x}$ where $x = x_{1} * \cdots * x_{n}$ (i.e. $x$ is a compound ideograph)and $s=a_{x} \sim x_{1} * \cdots * x_{n}$.
\end{definition}

The first reaction I felt is that there's something eerily self-referential here in the mechanism of the ateji - something very Godelian. But let us not get too excited yet.
